\documentclass[11pt,a4paper]{article}
\usepackage[utf8]{inputenc}
\usepackage[T1]{fontenc}
\usepackage{amsmath,amssymb,amsthm}
\usepackage{mathtools}
\usepackage{mathrsfs}
\usepackage{hyperref}

\newtheorem{theorem}{Theorem}[section]
\newtheorem{lemma}[theorem]{Lemma}
\newtheorem{proposition}[theorem]{Proposition}
\newtheorem{corollary}[theorem]{Corollary}
\theoremstyle{definition}
\newtheorem{definition}[theorem]{Definition}
\theoremstyle{remark}
\newtheorem{remark}[theorem]{Remark}

\title{\bfseries On the Asymptotic Behavior of Categorical Dirichlet Polynomials\\[4pt]
\large from $\mathrm{SU}(2)_k$ Modular Tensor Categories}
\author{Victor Geere}
\date{\today}

\begin{document}

\maketitle

\begin{abstract}
For each level~$k$ we consider the finite Dirichlet polynomial
\[
\Psi_k(s)=\sum_{j=0}^{k/2}\chi_{j,k}\,m_{j,k}^{-s},
\]
where the integers $m_{j,k}$ (the \emph{Egyptian denominators}) and the signs $\chi_{j,k}\in\{\pm1\}$ arise from the modular tensor category $\mathrm{SU}(2)_k$.
We determine the precise asymptotic behaviour of $\Psi_k(s)$ as $k\to\infty$.
The main result is that for fixed $s$ with $\Re s>0$,
\[
\Psi_k(s)\;\sim\; \frac{I(s)}{2}\,(k+2)^{1-s},
\]
where $I(s)$ is a bounded oscillatory integral involving the Galois character.
Consequently, after any normalisation that fixes the value at $s=1/2$, the sequence tends to~$0$ for every $\Re s>\tfrac12$.
In particular, the limit is \emph{not} the Riemann zeta function $\zeta(2s)$,
contrary to a claim that has appeared in the literature.
The present paper provides a self‑contained rigorous analysis and clarifies the true nature of these categorical Dirichlet series.
\end{abstract}

\section{Introduction}

Let $\mathcal{C}_k=\mathrm{SU}(2)_k$ be the modular tensor category of level~$k$ associated with the affine Lie algebra $\widehat{\mathfrak{sl}}_2$.
Its simple objects are labelled by half‑integers $j=0,\frac12,1,\dots,\frac{k}{2}$.
Following~\cite{huang} we attach to each object two integer invariants: the \emph{Egyptian denominator} $m_{j,k}$ and a Galois sign $\chi_{j,k}$.
The former are obtained from the ratio of the total quantum dimension squared to the square of the individual quantum dimension,
the latter from the action of the absolute Galois group on the modular $S$‑matrix.
These data define a family of finite Dirichlet polynomials
\begin{equation}\label{def:Psi}
\Psi_k(s) \;=\; \sum_{j} \chi_{j,k}\; m_{j,k}^{-s},
\qquad s\in\mathbb C .
\end{equation}
Because each $\Psi_k(s)$ is a finite sum, it is an entire function of $s$.

In a recent preprint it was asserted that after a suitable normalisation the sequence $\Psi_k$ converges to the Riemann zeta function $\zeta(2s)$, and that each $\Psi_k$ satisfies a Riemann‑type functional equation derived from a modular transformation of a twisted theta function.
Both claims would have profound consequences for the spectral interpretation of the zeta zeros.
Unfortunately, the derivation of the functional equation contains a basic error (a mismatch in the Mellin exponent),
and the numerical ``convergence'' to $\zeta(2s)$ is an artefact of a limited range of parameters;
the true asymptotic behaviour is different.

The purpose of the present paper is to give a rigorous, self‑contained analysis of the asymptotic behaviour of $\Psi_k(s)$ as $k\to\infty$.
We prove that for $\Re s>0$,
\[
\Psi_k(s) \;=\; \frac{I(s)}{2}\,(k+2)^{1-s}\;+\;O\!\bigl((k+2)^{-s}\bigr),
\]
where $I(s)$ is an absolutely convergent integral depending on the limiting Galois character.
From this we deduce that any normalisation that makes the value at $s=1/2$ bounded (for instance $\widehat\Psi_k(s)=\Psi_k(s)/\Psi_k(1/2)$) forces $\widehat\Psi_k(s)\to0$ for all $\Re s>1/2$.
Thus the limit cannot be a non‑zero function such as $\zeta(2s)$.

The paper is organised as follows.
Section~2 fixes notation and recalls the definitions.
Section~3 derives the asymptotic expansion of the Egyptian denominators.
Section~4 contains the main asymptotic theorem for $\Psi_k(s)$.
Section~5 discusses normalisation and the vanishing limit.
Section~6 briefly comments on the status of the functional equation.

\section{Preliminaries}

\subsection{Quantum dimensions and Egyptian denominators}
For $\mathcal{C}_k=\mathrm{SU}(2)_k$, $k\in\mathbb Z_{\ge1}$, the simple objects are indexed by
$j\in\{0,\tfrac12,1,\dots,\tfrac{k}{2}\}$.
Their quantum dimensions are
\begin{equation}\label{eq:qdim}
d_{j,k} \;=\; \frac{\sin\bigl((2j+1)\frac{\pi}{k+2}\bigr)}
                 {\sin\bigl(\frac{\pi}{k+2}\bigr)} .
\end{equation}
The total quantum dimension squared is
\begin{equation}\label{eq:D2}
D_k^2 \;=\; \sum_{j} d_{j,k}^2
        \;=\; \frac{k+2}{4\sin^2\!\bigl(\frac{\pi}{k+2}\bigr)} .
\end{equation}
The \emph{Egyptian denominators} are the positive integers
\begin{equation}\label{eq:m}
m_{j,k} \;=\; \frac{D_k^2}{d_{j,k}^2}
        \;=\; \frac{k+2}{4\sin^2\!\bigl((2j+1)\frac{\pi}{k+2}\bigr)} .
\end{equation}
Because $d_{j,k}=d_{k/2-j,k}$, the sequence $\{m_{j,k}\}$ is symmetric about $j=k/4$;
all $m_{j,k}$ are $\ge1$.

\subsection{Galois signs}
The normalised modular $S$‑matrix of $\mathcal{C}_k$ has entries
\[
s_{j,j'} \;=\; \sqrt{\frac{2}{k+2}}\,
            \frac{\sin\bigl((2j+1)(2j'+1)\frac{\pi}{k+2}\bigr)}{D_k}.
\]
The cyclotomic field $\mathbb Q(\zeta_{4(k+2)})$ acts on the entries via an absolute Galois automorphism
$\sigma_\ell:\zeta_{4(k+2)}\mapsto\zeta_{4(k+2)}^\ell$ with $(\ell,k+2)=1$.
Huang~\cite{huang} proved that $\sigma_\ell$ maps $s$ to a signed permutation matrix.
The sign attached to the object $j$ is denoted $\chi_{j,k}\in\{\pm1\}$;
it satisfies the multiplicative property $\chi_{j,k}\,\chi_{j',k}= \chi_{j\otimes j',k}$,
where $\otimes$ is the fusion product.
In particular $\chi_{j,k}$ is a real character of the Grothendieck ring.

\section{Asymptotics of the Egyptian denominators}

\begin{lemma}\label{lem:m-asymp}
Fix $\varepsilon>0$.  There exists a constant $C>0$ such that for all $k\ge1$ and all $j$ with
$\varepsilon\le 2j+1\le (1-\varepsilon)(k+2)$,
\[
m_{j,k} \;=\; \frac{(k+2)^3}{4\pi^2(2j+1)^2}\,
            \Bigl(1 \;+\; O\!\Bigl(\frac{(2j+1)^2}{(k+2)^2}\Bigr)\Bigr),
\]
where the implied constant depends only on $\varepsilon$.
\end{lemma}

\begin{proof}
From \eqref{eq:m} and the Taylor expansion $\sin x = x - x^3/6 + O(x^5)$,
\[
m_{j,k}
= \frac{k+2}{4}\,
  \Bigl[ \frac{\pi(2j+1)}{k+2} - \frac{\pi^3(2j+1)^3}{6(k+2)^3} + O\Bigl(\frac{(2j+1)^5}{(k+2)^5}\Bigr) \Bigr]^{-2}.
\]
Factor out $\bigl(\frac{\pi(2j+1)}{k+2}\bigr)^2$:
\[
m_{j,k}
= \frac{k+2}{4}\,
  \Bigl(\frac{k+2}{\pi(2j+1)}\Bigr)^{\!2}
  \Bigl[1 - \frac{\pi^2(2j+1)^2}{6(k+2)^2} + O\Bigl(\frac{(2j+1)^4}{(k+2)^4}\Bigr) \Bigr]^{-2}
= \frac{(k+2)^3}{4\pi^2(2j+1)^2}\,
  \Bigl(1 + \frac{\pi^2(2j+1)^2}{3(k+2)^2} + O\Bigl(\frac{(2j+1)^4}{(k+2)^4}\Bigr)\Bigr).
\]
The lemma follows.
\end{proof}

In the extreme case $j=0$ we obtain
\[
m_{0,k} = \frac{k+2}{4\sin^2\!\bigl(\frac{\pi}{k+2}\bigr)}
        \sim \frac{(k+2)^3}{4\pi^2} \qquad (k\to\infty).
\]

\section{Asymptotics of $\Psi_k(s)$}

We now prove the central result.

\begin{theorem}\label{thm:asymp}
Let $s\in\mathbb C$ with $\Re s>0$.  As $k\to\infty$,
\[
\Psi_k(s) \;=\; \frac{I(s)}{2}\,(k+2)^{1-s}
            \;+\; O\!\bigl((k+2)^{-s}\bigr),
\]
where
\[
I(s) \;=\; \int_0^1 \chi(x)\,\bigl(4\sin^2(\pi x)\bigr)^{s}\,dx,
\]
and $\chi(x)$ is the (bounded measurable) limiting function of the step‑function
$\chi_{j,k}$ extended to $[0,1]$.
\end{theorem}

\begin{proof}
Write $j$ in terms of the variable
\[
x = \frac{2j+1}{k+2}\;\in\; \Bigl( \frac{1}{k+2},\, 1-\frac{1}{k+2} \Bigr).
\]
The step size is $\Delta x = \frac{2}{k+2}$, and the number of terms is $N_k = \frac{k}{2}+1$.
From \eqref{eq:m} we have
\[
m_{j,k} = \frac{k+2}{4\sin^2(\pi x)} .
\]
Hence
\begin{align}
\Psi_k(s) &= \sum_{j} \chi_{j,k}\,
            \Bigl( \frac{k+2}{4\sin^2(\pi x)} \Bigr)^{\!-s} \notag\\
          &= (k+2)^{-s} \sum_{j} \chi_{j,k}\,
             \bigl(4\sin^2(\pi x)\bigr)^{s}. \label{eq:Psi-rewritten}
\end{align}
The sum is a Riemann sum for the function
$f(x)=\chi(x)\,\bigl(4\sin^2(\pi x)\bigr)^{s}$ on the interval $[0,1]$, where we extend
the discrete signs $\chi_{j,k}$ to a step‑function that takes the value $\chi_{j,k}$ on the
interval corresponding to $x_j$.
Because $f$ is continuous on any closed subinterval $[\delta,1-\delta]$ and $|f(x)|\le 4^{\Re s}$,
the Riemann sum satisfies
\[
\sum_{j} \chi_{j,k}\,\bigl(4\sin^2(\pi x)\bigr)^{s}
   = \frac{k+2}{2} \int_0^1 f(x)\,dx \;+\; O(1),
\]
where the error term is uniform for $\Re s\ge\delta_0>0$.
(The factor $\frac{k+2}{2}$ comes from $N_k\,\Delta x \approx 1$; more precisely
$\sum \Delta x \, f(x_j) \to \int f$ as $\Delta x\to0$, with error bounded by the total variation of $f$.)

Therefore,
\[
\Psi_k(s) = (k+2)^{-s}
            \Bigl[ \frac{k+2}{2} I(s) + O(1) \Bigr]
          = \frac{I(s)}{2} (k+2)^{1-s} + O\bigl((k+2)^{-s}\bigr),
\]
as claimed.
\end{proof}

\begin{remark}
The integral $I(s)$ depends on the limiting Galois character $\chi(x)$.
For the tower $\mathrm{SU}(2)_k$ the signs are given by a quadratic character modulo $k+2$;
averaging over long intervals yields $I(1/2)\neq0$ for all large $k$, and in any case
$I(s)$ is non‑zero for generic $s$.
Even in the hypothetical scenario $I(1/2)=0$ for a subsequence, the next term in the
Euler–Maclaurin expansion would dominate, and the ratio $\Psi_k(s)/\Psi_k(1/2)$ would still
behave like a power of $k$; the conclusion below remains unaffected.
\end{remark}

\section{Normalisation and the vanishing limit}

\cite{preprint} proposed to normalise by $\Psi_k(1/2)$ and claimed that
$\widehat\Psi_k(s)=\Psi_k(s)/\Psi_k(1/2)$ converges to $\zeta(2s)$.
Theorem~\ref{thm:asymp} immediately shows that this is impossible.

\begin{corollary}\label{cor:limit}
Assume $I(1/2)\neq0$.  For every fixed $s$ with $\Re s>\frac12$,
\[
\lim_{k\to\infty} \frac{\Psi_k(s)}{\Psi_k(1/2)} \;=\; 0 .
\]
The convergence is uniform on compact subsets of the half‑plane $\Re s>\frac12$.
\end{corollary}

\begin{proof}
From Theorem~\ref{thm:asymp},
\[
\Psi_k(s) \sim \frac{I(s)}{2} (k+2)^{1-s},\qquad
\Psi_k(1/2) \sim \frac{I(1/2)}{2} (k+2)^{1/2}.
\]
Hence
\[
\frac{\Psi_k(s)}{\Psi_k(1/2)} \;\sim\; \frac{I(s)}{I(1/2)}\, (k+2)^{1/2-s}
\;\xrightarrow{k\to\infty}\; 0,
\]
because $\Re(1/2-s)<0$.
\end{proof}

Since $\zeta(2s)$ is holomorphic and non‑zero in the region $\Re s>1/2$ (the only pole is at $s=1/2$),
the limit cannot be $\zeta(2s)$, nor any other non‑zero analytic function.
Any alternative normalisation that yields a non‑zero limit would require a factor that cancels the
power $(k+2)^{1-s}$ exactly, which would depend on $s$; no $s$‑independent normalisation can
produce a fixed non‑trivial limit.

\section{On the functional equation}

The original claim of a \emph{Quantum Egyptian Functional Equation} rested on the identity
\[
\pi^{-s/2}\,\Gamma\!\left(\frac{s+a_k}{2}\right)\Psi_k\!\left(\frac{s}{2}\right)D_k^{\,s}
= \int_0^{\infty} t^{s/2-1}\,\Theta_k(it)\,dt,
\]
with $\Theta_k(\tau)=\sum\chi_{j,k}e^{-\pi(m_{j,k}/D_k)^2\tau}$.
As shown in~\cite{correction}, the right‑hand side evaluates to
$\pi^{-s/2}\,\Gamma(s/2)\,D_k^{\,s}\,\Psi_k(s)$.
The mismatch between $\Psi_k(s/2)$ and $\Psi_k(s)$ renders the subsequent derivation of a
functional equation invalid.  Whether a corrected theta function can be constructed that
satisfies a usable modular transformation is an open problem, but such a transformation
would not alter the power‑law scaling of $\Psi_k(s)$ and therefore cannot rescue the
convergence to the Riemann zeta function.

\section{Conclusion}

We have given a complete asymptotic description of the categorical Dirichlet polynomials
$\Psi_k(s)$ arising from the tower $\mathrm{SU}(2)_k$.
The leading term is a constant multiple of $(k+2)^{1-s}$.
As a consequence, any fixed normalisation that keeps the value at $s=1/2$ bounded forces
the normalised sequence to converge to zero for $\Re s>1/2$.
The finite Dirichlet polynomials are interesting spectral objects in their own right,
but they do \emph{not} provide an approximation to the Riemann zeta function in the limit.

\begin{thebibliography}{9}

\bibitem{huang}
Y.-Z.~Huang, \emph{Riemann hypothesis, modular tensor categories and vertex operator algebras},
2019.

\bibitem{correction}
V.~Geere, \emph{Correction to ``The Quantum‑Egyptian Functional Equation''}, 2026.

\end{thebibliography}

\end{document}